\section{Task-Information Audit}
\label{sec:method}


\subsection{From wiring to the executed operator}
\label{sec:operator}

Let $A_{uv}=1$ denote a directed edge $u\!\to\!v$. The main experiments
use binary support and incoming-neighbor averaging,
\begin{equation}
 d_v=\sum_u A_{uv},\qquad
 W_{vu}=\frac{A_{uv}}{\max(d_v,1)}.
 \label{eq:operator}
\end{equation}
Rows for zero in-degree remain zero. Features propagate linearly for two
rounds, $H=W^2X$, before an MLP reads the concatenated features at
$S=\{s_1,s_2\}$. The operator direction, normalization, and readout
locations are part of the computational model.

With nonzero features only at sources $a,b$, let $E_{ab}$ place the source
vectors and $R_S$ select the sinks. The graph affects the classifier input
through
\begin{equation}
 H_S=T(A)X_{ab},\qquad T(A)=R_SW(A)^2E_{ab}.
 \label{eq:transfer}
\end{equation}
$T$ has two source columns and two sink rows. It records every two-step
route contributing to these inputs, including routes outside a selected
motif family. Equal $T$ matrices give identical readout inputs for any
source features; equal graph statistics need not. Computing $T$ requires
propagating two basis columns, rather than constructing all entries of
$W^2$. For a sparse fixed operator and $h$ steps, this costs
$O(h|E|\,n_{\rm source})$ time and $O(Nn_{\rm source})$ storage.

\subsection{Source visibility and ideal task solvability}
\label{sec:observability}

A source is \emph{visible} when its column of $T$ is nonzero. This is readout
access for the fixed operator. It is not classical state observability, and it
does not recover the full source state. For the ideal noiseless task, source
visibility gives a complete criterion.
Let $c_a,c_b$ be independent balanced bits and
$x_a=P_{c_a}z_a$, $x_b=P_{c_b}z_b$, with independent
$z_a,z_b\sim\mathcal N(0,I_8)$. The class subspaces
$U_c=\operatorname{span}(P_c)\subset\mathbb R^{16}$ satisfy
$U_0\cap U_1=\{0\}$, and the label is $y=c_a\oplus c_b$.

\paragraph{Remark (common-relay sufficiency).}
If $\mathcal C(A)$ in Eq.~\eqref{eq:paths} is nonempty, both source columns
of a nonnegative $T$ are nonzero. The converse is false: separate relays can
expose both sources without any shared relay.

\paragraph{Proposition 1 (ideal population criterion).}
For a fixed nonnegative $T$ and no other active inputs, the Bayes accuracy
from $H_S=TX_{ab}$ is
\begin{equation}
 \operatorname{Acc}^{*}(T)=
 \begin{cases}
  1, & T_{:a}\ne0\ \text{and}\ T_{:b}\ne0,\\
  \tfrac12, & \text{otherwise}.
 \end{cases}
 \label{eq:criterion}
\end{equation}

\emph{Proof.} If $T_{:b}=0$, the readout is independent of $c_b$,
so $\Pr(y=1\mid H_S)=1/2$; the same holds for $a$.
If both columns are nonzero, summing sink vectors gives
$z=\alpha x_a+\beta x_b$, where
$\alpha=\sum_sT_{sa}>0$ and $\beta=\sum_sT_{sb}>0$.
For equal source classes, $z\in U_0\cup U_1$. For different classes,
transversality and nonzero Gaussian coordinates imply
$z\notin U_0\cup U_1$ almost surely. Subspace membership therefore
recovers the XOR label. $\square$

This is a population information statement. Finite MLP optimization,
sample size, numerical precision and noise can change achieved accuracy.
The supplement gives the full argument and the ten-seed basis-rank check.
A shared relay and rank-two transfer are both unnecessary: separate
source routes suffice, and even identical sink mixtures retain the ideal
label information when both source coefficients are positive.

Noise, finite readouts and numerical precision can change achieved accuracy;
the criterion isolates whether the executed operator supplies the two source
classes to the readout in the first place.

\subsection{What the structural interventions constrain}
\label{sec:paths}

The graph controls preserve each node's in-/out-degree and module-pair
edge counts. For fixed sources and sinks, a common-relay family is
\begin{equation}
 \mathcal C(A)=\{(r,s):s\in S,\ A_{ar}A_{br}A_{rs}=1\}.
 \label{eq:paths}
\end{equation}
A member represents two source edges converging at a relay and one edge
to a sink. Its three-edge description omits separate-relay routes.
The number of named motif edges retained is an even narrower statistic.
Neither quantity specifies all columns of Eq.~\eqref{eq:transfer}.

We compare two archived intervention designs. Named-edge compensation
removes requested edges of a selected four-edge motif and replaces each
through a module-matched directed double-edge swap,
\begin{equation}
 (u\!\to\!v,x\!\to\!y)
 \longmapsto (u\!\to\!y,x\!\to\!v).
 \label{eq:swap}
\end{equation}
Matching source-module and target-module labels preserves module counts;
the exchange preserves every labeled degree. It does not constrain
other source-specific routes.

The path-conditioned design enumerates $\mathcal C(A)$, retains its first
$q$ relay/sink pairs in a fixed order, and protects their path-edge union.
Edges belonging only to non-retained triples are processed in sorted order.
Background partners are randomized subject to degree/module matching,
no duplicate/self-loop edges, no direct source--sink edges, and protection
of the retained path edges. A candidate is additionally rejected if it
creates any new member of $\mathcal C$; the final set must exactly equal
the requested retained subset. Infeasible requests are recorded.
This is a finite common-relay constraint, not a complete source-transport
constraint. Different accepted swaps and ordering choices can still
change source-specific routes and hence $T$.

Subsequent degree nulls rewire the whole dose graph. Protected nulls fix
the retained path-edge union while rewiring the remainder. They are not
subject to the no-new-common-relay acceptance rule and need not preserve
$T$. Table~\ref{tab:questions} states the distinct computational questions
answered by these counterfactuals.

\begin{table}[t]
\centering
\caption{\textbf{Specify the computational contrast.} $A_q$ denotes a
path-conditioned dose graph. The task audit checks the executed transfer
alongside the declared graph constraints.}
\label{tab:questions}
\small
\begin{tabular}{@{}p{0.30\linewidth}p{0.64\linewidth}@{}}
\toprule
Comparison & Question answered \\
\midrule
$A$ versus $W=I$ & Is cross-node transport required at the chosen readout? \\[3pt]
$A$ versus degree null & Does this graph outperform the specified rewiring? \\[3pt]
$A_q$ versus $A_{q'}$ & What changes under this particular constrained edit? \\[3pt]
$A_q$ versus protected null & What changes when selected edges stay fixed but the background is rewired? \\
\bottomrule
\end{tabular}
\end{table}

\subsection{Reading an intervention}
\label{sec:controls}

The audit has a direct sequence. First establish the executed direction,
normalization, propagation depth, and readout. Next enumerate the stated
motif family and compute source-specific transfer to those readout nodes.
Finally distinguish loss of source access from the performance of a
particular trained readout. An oracle given both sources tests label
learnability; $W=I$ prevents cross-node transport to distinct sinks;
linear-readout, shuffled-label, and moved-readout controls check alternative
explanations. Equation~\eqref{eq:criterion} then separates ideal label-information
loss from the trainability of the surviving input mixture.

The exact matrix reduction applies to a fixed linear propagation prefix.
Nonlinear message passing, gating and spiking dynamics require an
input-dependent analysis of source access.
